\documentclass{beamer}\usepackage[]{graphicx}\usepackage[]{color}
%% maxwidth is the original width if it is less than linewidth
%% otherwise use linewidth (to make sure the graphics do not exceed the margin)
\makeatletter
\def\maxwidth{ %
  \ifdim\Gin@nat@width>\linewidth
    \linewidth
  \else
    \Gin@nat@width
  \fi
}
\makeatother

\definecolor{fgcolor}{rgb}{1, 0.894, 0.769}
\newcommand{\hlnum}[1]{\textcolor[rgb]{0.824,0.412,0.118}{#1}}%
\newcommand{\hlstr}[1]{\textcolor[rgb]{1,0.894,0.71}{#1}}%
\newcommand{\hlcom}[1]{\textcolor[rgb]{0.824,0.706,0.549}{#1}}%
\newcommand{\hlopt}[1]{\textcolor[rgb]{1,0.894,0.769}{#1}}%
\newcommand{\hlstd}[1]{\textcolor[rgb]{1,0.894,0.769}{#1}}%
\newcommand{\hlkwa}[1]{\textcolor[rgb]{0.941,0.902,0.549}{#1}}%
\newcommand{\hlkwb}[1]{\textcolor[rgb]{0.804,0.776,0.451}{#1}}%
\newcommand{\hlkwc}[1]{\textcolor[rgb]{0.78,0.941,0.545}{#1}}%
\newcommand{\hlkwd}[1]{\textcolor[rgb]{1,0.78,0.769}{#1}}%
\let\hlipl\hlkwb

\usepackage{framed}
\makeatletter
\newenvironment{kframe}{%
 \def\at@end@of@kframe{}%
 \ifinner\ifhmode%
  \def\at@end@of@kframe{\end{minipage}}%
  \begin{minipage}{\columnwidth}%
 \fi\fi%
 \def\FrameCommand##1{\hskip\@totalleftmargin \hskip-\fboxsep
 \colorbox{shadecolor}{##1}\hskip-\fboxsep
     % There is no \\@totalrightmargin, so:
     \hskip-\linewidth \hskip-\@totalleftmargin \hskip\columnwidth}%
 \MakeFramed {\advance\hsize-\width
   \@totalleftmargin\z@ \linewidth\hsize
   \@setminipage}}%
 {\par\unskip\endMakeFramed%
 \at@end@of@kframe}
\makeatother

\definecolor{shadecolor}{rgb}{.97, .97, .97}
\definecolor{messagecolor}{rgb}{0, 0, 0}
\definecolor{warningcolor}{rgb}{1, 0, 1}
\definecolor{errorcolor}{rgb}{1, 0, 0}
\newenvironment{knitrout}{}{} % an empty environment to be redefined in TeX

\usepackage{alltt}
\usepackage{../371g-slides}
\title{Decision Trees 1}
\subtitle{Lecture 6}
\author{STA 371G}
\usepackage{forest}
\IfFileExists{upquote.sty}{\usepackage{upquote}}{}
\begin{document}
  
  

  \frame{\maketitle}

  % Show outline at beginning of each section
  \AtBeginSection[]{
    \begin{frame}<beamer>
      \tableofcontents[currentsection]
    \end{frame}
  }

  %%%%%%% Slides start here %%%%%%%

  \begin{darkframes}

    %{slide 1}
    \begin{frame}{Making Better Decisions}
      \fontsize{10}{10}\selectfont
      \begin{center}
        \includegraphics[width=2in]{DecisionAnalysis.png} \\
      \end{center}
        \textit{Decision making is the only way that individuals can purposely
        exercise any control over their lives, careers, or their surroundings.}
        \text{- Ralph Keeney, Making Better Decision Makers, Decision Analysis, vol. 1}
        No:4, 2004

      \lc %{How do you make important decisions?}
    \end{frame}


    %{slide 2}
    \begin{frame}[fragile]{Decision Analysis}
      \begin{itemize}[<+->]
        \item A framework for analyszing decision problems that involve uncertainty
        \item Decision trees: a powerful graphical tool that guides that analysis
        \item Smaller analyses can be done using pen and paper
        \item Larger ones require software
        \end{itemize}
    \end{frame}


    %{slide 3}
    \begin{frame}[fragile]{Topics to cover}
          \begin{itemize}[<+->]
            \item Criteria for choosing among alternative decisions
            \item How probabilities are used in the decision-making process
            \item How early decisions affect later decisions
            \item How a decision-maker can quantify the value of information
            \item How attitudes toward risk and uncertainty can affect the analysis
          \end{itemize}

    \end{frame}


    %{slide 4}
    \begin{frame}[fragile]{Elements of a Decision Analysis}
      All problems have three common elements:

        \begin{itemize}
            \item The decisions available to the decision maker.
            \item The possible outcomes and the probabilities of these outcomes.
            \item A value model that provides monetary values for the various outcomes.
        \end{itemize}

      Once these elements are defined, the decision maker can find an optimal decision.
    \end{frame}


    %{slide #5}
    \begin{frame}[fragile]{Payoff Tables}
      A payoff table lists the payoff for each decision outcome pair; positive values are gains and negative values are losses.

        \begin{center}
          \begin{tabular}{r|rrr}
            & O1 & O2 & O3 \\
            \hline
            D1 & \$10 & \$10 & \$10 \\
            D2 & $-\$10$ & \$20 & \$30 \\
            D3 & $-\$30$ & \$30 & \$80
          \end{tabular}
        \end{center}
        \begin{itemize}
            \item This table shows three possible decisions (D1, D2, and D3) and three
            possible outcomes (O1, O2, and O3) for each.
            \item Which decision do you prefer?
        \end{itemize}

      \lc %{Which decision do you prefer?}
    \end{frame}


    %{slide #6}
    \begin{frame}[fragile]{Payoff Tables}
      We need to know the probability of each outcome to make a good decision!
      \begin{center}
        \begin{tabular}{r|rrr}
          & O1 & O2 & O3 \\
          \hline
          D1 & \$10 & \$10 & \$10 \\
          D2 & $-\$10$ & \$20 & \$30 \\
          D3 & $-\$30$ & \$30 & \$80
        \end{tabular}
      \end{center}
        \begin{itemize}
            \item Suppose $P(O1) = 0.3, P(O2) = 0.5, P(O3) = 0.2$
            \item Now which decision do you prefer?
        \end{itemize}

    \end{frame}


    %{slide 7}
    \begin{frame}[fragile]{Expected Value}
         \begin{itemize}[<+->]
            \item Expected value (EV) is a weighted average of the possible payoffs for the decision,
            weighted by the probabilities of the outcomes
            \item $E(D1) = 10$
            \item $E(D2) = -10(0.3) + 20(0.5) - 30(0.2) = 13$
            \item $E(D3) = ?$
        \end{itemize}

      \lc %{What is EV of D3?}
    \end{frame}


    %{slide 8}
    \begin{frame}[fragile]{Decision Trees}
          \begin{itemize}
            \item Time proceeds from left to right. \pause
            \item Branches leading out of a decision node represent the possible decisions. \pause
            \item Probabilities are listed on probability branches, and are conditional on the events that have already been observed (i.e., they assume that everything to the left has already happened). \pause
            \item Monetary values are shown to the right of the end nodes. \pause
            \item EVs are calculated through a ``rolling-back'' process.
          \end{itemize}
    \end{frame}


    %{slide 9}
    \begin{frame}[fragile]{Example}
    \begin{center}
        \begin{tikzpicture}
        [
          grow                    = right,
          level 1/.style          ={sibling distance=4em, level distance=8em},
          level 2/.style          ={sibling distance=3em, level distance=14em},
          level 3/.style          ={sibling distance=3em, level distance=8em},
          %level distance          = 8em,
          edge from parent/.style = {draw, -latex},
          every node/.style       = {font=\footnotesize},
          sloped
        ]
        \node [decision] {}
          child {
            node [prob] {}
            child {
              node [leaf] { $80$ }
              edge from parent node [above] {O3, $p=0.2$}
            }
            child {
              node [leaf] { $30$ }
              edge from parent node [above] {O2, $p=0.5$}
            }
            child {
              node [leaf] { $-30$ }
              edge from parent node [above] {O1, $p=0.3$}
            }
            edge from parent node [above] {D3}
          }
          child {
            node [leaf] {$10$}
            edge from parent node [above] {D1}
          }
          child {
            node [prob] {}
            child {
              node [leaf] { $30$ }
              edge from parent node [above] {O3, $p=0.2$}
            }
            child {
              node [leaf] { $20$ }
              edge from parent node [above] {O2, $p=0.5$}
            }
            child {
              node [leaf] { $-10$ }
              edge from parent node [above] {O1, $p=0.3$}
            }
            edge from parent node [above] {D2}
          };
        \end{tikzpicture}
      \end{center}

    \end{frame}


        %{slide 10}
    \begin{frame}[fragile]{Rolling back: Step 1}
    Calculate the expected value at each probability node:
    \begin{center}
        \begin{tikzpicture}
        [
          grow                    = right,
          level 1/.style          ={sibling distance=4em, level distance=8em},
          level 2/.style          ={sibling distance=3em, level distance=14em},
          level 3/.style          ={sibling distance=3em, level distance=8em},
          %level distance          = 8em,
          edge from parent/.style = {draw, -latex},
          every node/.style       = {font=\footnotesize},
          sloped
        ]
        \node [decision] {}
          child {
            node [prob] {?}
            child {
              node [leaf] { $80$ }
              edge from parent node [above] {O3, $p=0.2$}
            }
            child {
              node [leaf] { $30$ }
              edge from parent node [above] {O2, $p=0.5$}
            }
            child {
              node [leaf] { $-30$ }
              edge from parent node [above] {O1, $p=0.3$}
            }
            edge from parent node [above] {D3}
          }
          child {
            node [leaf] {$10$}
            edge from parent node [above] {D1}
          }
          child {
            node [prob] {13}
            child {
              node [leaf] { $30$ }
              edge from parent node [above] {O3, $p=0.2$}
            }
            child {
              node [leaf] { $20$ }
              edge from parent node [above] {O2, $p=0.5$}
            }
            child {
              node [leaf] { $-10$ }
              edge from parent node [above] {O1, $p=0.3$}
            }
            edge from parent node [above] {D2}
          };
        \end{tikzpicture}
      \end{center}
      $E(D1) = .3(-10) + .5(20) + .2(30) = 13$
      \lc
    \end{frame}



            %{slide 11}
    \begin{frame}[fragile]{Rolling back: Step 2}
     Calculate the maximum at each decision node:
    \begin{center}
        \begin{tikzpicture}
        [
          grow                    = right,
          level 1/.style          ={sibling distance=4em, level distance=8em},
          level 2/.style          ={sibling distance=3em, level distance=14em},
          level 3/.style          ={sibling distance=3em, level distance=8em},
          %level distance          = 8em,
          edge from parent/.style = {draw, -latex},
          every node/.style       = {font=\footnotesize},
          sloped
        ]
        \node [decision] {22}
          child {
            node [prob] {22}
            child {
              node [leaf] { $80$ }
              edge from parent node [above] {O3, $p=0.2$}
            }
            child {
              node [leaf] { $30$ }
              edge from parent node [above] {O2, $p=0.5$}
            }
            child {
              node [leaf] { $-30$ }
              edge from parent node [above] {O1, $p=0.3$}
            }
            edge from parent node [above] {D3}
          }
          child {
            node [leaf] {$10$}
            edge from parent node [above] {D1}
          }
          child {
            node [prob] {13}
            child {
              node [leaf] { $30$ }
              edge from parent node [above] {O3, $p=0.2$}
            }
            child {
              node [leaf] { $20$ }
              edge from parent node [above] {O2, $p=0.5$}
            }
            child {
              node [leaf] { $-10$ }
              edge from parent node [above] {O1, $p=0.3$}
            }
            edge from parent node [above] {D2}
          };
        \end{tikzpicture}
      \end{center}
       Take decision D3 since $22 = \text{max}(10, 13, 22)$.
    \end{frame}







    %{slide 12}
    \begin{frame}[fragile]{Sally Ann Soles' Shoe Factory}
      Sally Ann Soles manages a successful shoe factory. She is wondering
      whether to expand her factory this year.
      \begin{itemize}
        \item The cost of the expansion is \$1.5M. \pause

        \item If she does nothing and the economy stays good, she expects to earn \$3M in
        revenue, but if the economy is bad, she expects only \$1M. \pause

        \item If she expands the factory, she expects to earn \$6M if the economy is
        good and \$2M if it is bad. \pause

        \item She estimates that there is a 40 percent chance of a good economy and
        a 60 percent chance of a bad economy.
      \end{itemize}
     Should she expand?
    \end{frame}


    %{slide 13}
    \begin{frame}
      \begin{center}
        \begin{tikzpicture}
        [
          grow                    = right,
          level 1/.style          ={sibling distance=8em},
          level 2/.style          ={sibling distance=6em},
          level 3/.style          ={sibling distance=6em},
          level distance          = 12em,
          edge from parent/.style = {draw, -latex},
          every node/.style       = {font=\footnotesize},
          sloped
        ]
        \node [decision] {}
          child {
              node [prob] {}
              child{
                node [leaf] {\$1M}
                 edge from parent node [above] {Bad Economy}
                 edge from parent node [below] {$p=0.6$}
              }
              child{
                node [leaf] {\$3M}
                 edge from parent node [above] {Good Economy}
                 edge from parent node [below] {$p=0.4$}
              }
                edge from parent node [above] {Don't Expand}
                edge from parent node [below] {spend \$0}
          }
          child {
            node [prob] {}
            child{
              node [leaf] {\$0.5M}
               edge from parent node [above] {Bad Economy}
               edge from parent node [below] {$p=0.6$}
            }
            child{
              node [leaf] {\$4.5M}
               edge from parent node [above] {Good Economy}
               edge from parent node [below] {$p=0.4$}
            }
            edge from parent node [above] {Expand Factory}
            edge from parent node [below] {spend \$1.5M}
          };
        \end{tikzpicture}
      \end{center}
      \pause
      $E(\text{expand}) = .4(4.5) + .6(0.5) = 2.1$
      $E(\text{don't expand}) = .4(3) + .6(1) = 1.8$
    \end{frame}



      %{slide 13-2}
          \begin{frame}
      \begin{center}
        \begin{tikzpicture}
        [
          grow                    = right,
          level 1/.style          ={sibling distance=8em},
          level 2/.style          ={sibling distance=6em},
          level 3/.style          ={sibling distance=6em},
          level distance          = 12em,
          edge from parent/.style = {draw, -latex},
          every node/.style       = {font=\footnotesize},
          sloped
        ]
        \node [decision] {\$2.1M}
          child {
              node [prob] {\$1.8M}
              child{
                node [leaf] {\$1M}
                 edge from parent node [above] {Bad Economy}
                 edge from parent node [below] {$p=0.6$}
              }
              child{
                node [leaf] {\$3M}
                 edge from parent node [above] {Good Economy}
                 edge from parent node [below] {$p=0.4$}
              }
                edge from parent node [above] {Don't Expand}
                edge from parent node [below] {spend \$0}
          }
          child {
            node [prob] {\$2.1M}
            child{
              node [leaf] {\$0.5M}
               edge from parent node [above] {Bad Economy}
               edge from parent node [below] {$p=0.6$}
            }
            child{
              node [leaf] {\$4.5M}
               edge from parent node [above] {Good Economy}
               edge from parent node [below] {$p=0.4$}
            }
            edge from parent node [above] {Expand Factory}
            edge from parent node [below] {spend \$1.5M}
          };
        \end{tikzpicture}
      \end{center}
      $E(\text{expand}) = (.4(6) + .6(2)) - 1.5 = 2.1$
      $E(\text{don't expand}) = (.4(3) + .6(1)) = 1.8$

      Since \$2.1 > \$1.8, she should expand!
    \end{frame}



    %{slide 14}
    \begin{frame}[fragile]{Sequential decisions}
      She later learns if she expands, she can opt to either:
      \begin{enumerate}[(a)]
        \item expand the factory further, which costs \$1.5M and will yield an extra \$2M in profit if the economy is good, but \$1M if it is bad,   \pause
        \item abandon the project and sell the equipment she originally bought, for \$1.3M, or  \pause
        \item do nothing.
      \end{enumerate}
      How has the decision changed?
    \end{frame}




    \begin{frame}
      \begin{center}
        \begin{tikzpicture}
        [
          grow                    = right,
          level 1/.style          ={sibling distance=8em, level distance=5em},
          level 2/.style          ={sibling distance=11em, level distance=8em},
          level 3/.style          ={sibling distance=5em, level distance=12em},
          %level distance          = 8em,
          edge from parent/.style = {draw, -latex},
          every node/.style       = {font=\footnotesize},
          sloped
        ]
        \node [decision] {}
            child {
              node [prob] {\$1.8M}
              edge from parent node [above] {Don't Expand}
              edge from parent node [below] {spend \$0}
            }
            child{
              node [prob] { }
              child{
                node [decision] { }
                child{
                  node [leaf] {\$0.8M}
                  edge from parent node [above] {Reduce \& Sell}
                  edge from parent node [below] {+\$1.3M}
                }
                child{
                  node [leaf] {\$0.5M}
                  edge from parent node [above] {Stay at new level}
                }
                child{
                  node [leaf] {\$0}
                  edge from parent node [above] {Expand further}
                  edge from parent node [below] {-\$1.5M}
                }
                 edge from parent node [above] {Bad Economy}
                 edge from parent node [below] {$p=0.6$}
              }
              child{
                node [decision] { }
                child{
                  node [leaf] {\$2.8M}
                  edge from parent node [above] {Reduce \& Sell}
                  edge from parent node [below] {+\$1.3M}
                }
                child{
                  node [leaf] {\$4.5M}
                  edge from parent node [above] {Stay at new level}
                }
                child{
                  node [leaf] {\$5M}
                  edge from parent node [above] {Expand further}
                  edge from parent node [below] {-\$1.5M}
                }
                 edge from parent node [above] {Good Economy}
                 edge from parent node [below] {$p=0.4$}
              }
              edge from parent node [above] {Expand}
              edge from parent node [below] {-\$1.5M}
            };
        \end{tikzpicture}
      \end{center}
      \note{Build this up together on the board so students}
      \note{see where each value comes from. The payoffs}
      \note{are tricky, esp for the reduce & sell cases.}
      \note{Use the last LC question one or more times}
      \note{to prompt students.}
    \end{frame}





    %{slide 16}
    \begin{frame}[fragile]{Expected value of the option}
        The EV of expanding is now
        \[ (.4(5) + .6(0.8)) = 2.48. \]
        \pause
        Therefore the value of the option is
        \[ E(\text{new}) - E(\text{old}) = 2.48 - 2.1 = 0.38, \]
        or \$380,000.
    \end{frame}


   %{slide 17}
    \begin{frame}[fragile]{Does this look familiar?}
      \begin{itemize}[<+->]
        \item This method of valuing a real option is used by corporate strategists, management consultants, and bankers.
        \item A further refinement is to take time value of money into account and present value everything.
      \end{itemize}
    \end{frame}
  \end{darkframes}
\end{document}
